\selectlanguage{english}
\chapter*{Abstract}
\addcontentsline{toc}{chapter}{Abstract}

This research develops and validates a methodology based on technical and
regulatory criteria for selecting user-interface technologies within the
context of digital transformation in banking institutions, focusing on a
comparative analysis of Angular, React, and Vue.js. In a sector where
security, performance, and regulatory compliance are critical, technology
selection represents a strategic decision with direct impact on
competitiveness and regulatory exposure.

An additional factor motivates and delimits the scope of this research: the
consolidation of AI-assisted code-generation tools (GitHub Copilot, Cursor,
and equivalent assistants), which progressively narrows the productivity and
learning-curve gap traditionally associated with developer talent
availability --an organizational factor-- shifting the decisive axis of
technology selection toward the framework's intrinsic technical properties
and its regulatory fit, neither of which AI can substitute. Accordingly, the
proposed model is structured around two dimensions --technical and
regulatory-- comprising ten evaluation criteria, two of them emergent: the
quality and auditability of AI-generated code (technical dimension) and the
regulatory governance of generative-AI tool usage in banking software
development (regulatory dimension).

The proposed methodology consists of four integrated phases: institutional
requirements analysis, technical evaluation (including AI-assisted code
quality), regulatory and compliance evaluation, and a structured final
decision-making process. To validate this approach, a case study was
conducted in a mid-sized banking institution undergoing modernization of its
digital channels, systematically applying the criteria and tools defined in
the methodology, and contrasting results with and without the criteria
associated with generative AI.

The results show that, under the technical-regulatory model, Angular obtains
the highest weighted score ($S=8.84$), followed by React ($S=7.83$) and
Vue.js ($S=7.56$), owing to its typed, opinionated architecture, which favors
both code security and maintainability and the traceability required by
regulatory audits, while additionally producing more consistent and
verifiable AI-assisted code. The sensitivity analysis confirms this
preference is robust: unlike a model that also weighs organizational
factors --where an extreme scenario could favor React-- the proposed
technical-regulatory model keeps Angular as the optimal choice across every
plausible scenario for a regulated banking environment.

This research contributes to the field by providing a methodological
framework specifically adapted to the needs of the banking sector in the era
of generative AI, offering assessment and decision-making tools that can be
adopted by financial institutions facing similar digital transformation
challenges. It concludes that the emergence of AI code assistants does not
eliminate the need for a rigorous technology-selection process, but instead
reconfigures its decisive criteria toward the technical and regulatory axes.

\vspace{0.5cm}
\noindent \textbf{Keywords:} Digital banking transformation, user-interface
technologies, Angular, React, Vue, technical and regulatory criteria,
AI-assisted code generation, technology selection methodology, financial
applications.

\selectlanguage{spanish}
